\documentclass[10pt,a4paper]{amsart}
\usepackage[margin=19mm]{geometry}
\usepackage{amsmath,amssymb,mathtools}
\usepackage[scr=rsfs,cal=euler]{mathalfa}
\usepackage{xfrac}
\usepackage[unicode,hidelinks,bookmarks=false]{hyperref}
\newcommand{\ie}{i.e.\ }
\newcommand{\cf}{cf.\ }
\newcommand{\resp}{respectively\ }
\newcommand{\Subsection}{\subsection}
\providecommand{\nobreakdash}{\nobreak\hbox{-}}
\setlength{\emergencystretch}{2em}
\multlinegap0pt
\usepackage{lmodern}
\usepackage{mathtools}
\usepackage{enumitem,booktabs}
\newtheorem{theorem}{\bf Theorem}[section]
\newtheorem{proposition}[theorem]{\bf Proposition}
\newtheorem{lemma}[theorem]{\bf Lemma}
\newtheorem{corollary}[theorem]{\bf Corollary}
\newtheorem{conjecture}[theorem]{\bf Conjecture}
\newtheorem{problem}[theorem]{\bf Problem}
\newtheorem{example}[theorem]{\bf Example}
\newtheorem{examples}[theorem]{\bf Examples}
\newtheorem{definition}[theorem]{\bf Definition}
\newtheorem{remark}[theorem]{\bf Remark}
\newcommand{\dd}{\mathsf d}
\newcommand{\DD}{\mathsf D}
\newcommand{\height}{\mathsf h}
\DeclareMathOperator{\ord}{ord}
\DeclareMathOperator{\supp}{supp}
\DeclareMathOperator{\Aut}{Aut}
\begin{document}
\title[On a classical zero-sum invariant II: Disproof of a long-sta]{On a classical zero-sum invariant II: Disproof of a long-standing conjecture}
\author{Alfred Geroldinger \and Guoqing Wang \and Wenkai Yang}
\date{}
\maketitle
\noindent\emph{Source and licence.} On a classical zero-sum invariant II: Disproof of a long-standing conjecture. Version arXiv:2609.17127v1 (CC BY 4.0). This material is distributed under the Creative Commons Attribution 4.0 International licence, http://creativecommons.org/licenses/by/4.0/, as recorded in the arXiv OAI licence metadata for this exact version and shown on the abstract page https://arxiv.org/abs/2609.17127v1. Full rights notice: licenses/arxiv-2609.17127-CC-BY-4.0.txt.\par\smallskip
\noindent\textit{Editorial scope.} Counted unit: the original \texttt{theorem} environment \texttt{thm:main2} at source lines 109--111 together with its complete proof at lines 222--312; the delivered document keeps source lines 70--313 so that every referenced label and every citation resolves inside it. Native editable source is the paper's own concatenated TeX; the AMS wrapper is editorial formatting only. No publisher class, upstream script or unrelated artwork is redistributed.
\section{Introduction}


Let $G$ be a nontrivial finite abelian group, say $G = C_{n_1} \oplus \ldots \oplus C_{n_r}$, where $r=\mathsf r (G) \in \mathbb N$ is the rank of $G$ and $1 < n_1 \mid \ldots \mid n_r$. We set
$G^\bullet=G\setminus\{0\}$, $\mathsf d^* (G) = \sum_{i=1}^r (n_i-1)$, and  $\mathsf D^* (G) = 1 + \mathsf d^* (G)$. The (small) Davenport constant $\mathsf d (G)$ is the maximal length of a zero-sum free sequence over $G$,  the (large) Davenport constant $\mathsf D (G)$ is the maximal length of a minimal zero-sum sequence over $G$, and we have $\mathsf D (G) = \mathsf d (G)+1 \ge \mathsf D^* (G)$. It is well-known that $\mathsf D (G) = \mathsf D^* (G)$ if $r \le 2$, if $G$ is a $p$-group, and for many others. But, in general inequality may hold. If $r \le 2$, then the structure of minimal zero-sum sequences of maximal length is known. But, if $r \ge 3$, even for $p$-groups only little is known about the structure of such sequences.

The invariant $\nu (G)$ (introduced by Emde Boas in the 1960s, \cite{Em69a}) offers information on the structure of extremal sequences by considering the set of their subsequence sums. Let $p$ be a prime divisor of $|G|$. Let $\nu (G)$ resp. $\nu_p (G)$ be the smallest integer such that, for any zero-sum free sequence $T$ over $G$, $|T| \ge \nu (G)$ resp. $|T| \ge \nu_p (G)$ implies that
\begin{equation} \label{def-nu(G)}
G^{\bullet} \setminus \Sigma (T) \subseteq \alpha+H \ \mbox {for some  subgroup $H  \subsetneq G$ and some $\alpha\in G\setminus H$}
\end{equation}
resp.
\begin{equation} \label{def-nu_p(G)}
G^{\bullet} \setminus \Sigma (T) \subseteq \alpha + H \ \text{for some subgroup} \ H \subseteq G \ \text{with} \ [G \colon H] = p \ \text{and some} \ \alpha \in G \setminus H \,.
\end{equation}
If $T \in \mathcal F (G)$ is zero-sum free with $|T|=\mathsf d(G)$, then, by a straightforward argument, we get $\Sigma (T) = G^{\bullet}$, whence, by definition, $\nu (G) \le \nu_p (G) \le \mathsf d (G)$. A further simple argument shows that
\begin{equation} \label{basic-inequ-2}
\mathsf d (G) -1 \le \nu (G) \le \nu_p (G) \le \mathsf d (G) \,.
\end{equation}



In 2000 Weidong Gao proposed the following conjecture (see \cite[Section~3]{Ga00b} or \cite[Conjecture~4.10]{Ga-Ge06b}).
\begin{conjecture} \label{conj:gao}
Every nontrivial finite abelian group $G$ satisfies $\nu(G)=\dd(G)-1$.
\end{conjecture}


This conjecture initiated a great deal of  further research. It got confirmed for $p$-groups (\cite[Theorem 5.5.9]{Ge-HK06a}), for  groups of rank at most two (\cite[Theorem 3.4.11]{Ge-Gr-Zh26a}), and others. Only recently it was confirmed for a series of groups $G$ satisfying $\mathsf d (G) > \mathsf d^* (G)$ (\cite{Ge-Ya27a}).


In this note, we disprove the conjecture by providing a series of groups with $\nu (G) = \mathsf d (G)$.

\begin{theorem}\label{thm:main1}
Let $G=C_2^3\oplus C_{2n}$ for some odd $n\geq 3$. Then $$\nu(G)=\nu_2 (G)=\dd(G)=\mathsf d^* (G).$$
\end{theorem}

Note that, if $G = C_2^r \oplus C_{2n}$ with $n>70$ odd, then $\nu (G) = \mathsf d (G)-1$ for $r \in \{0,1,2,4\}$. For the growth of $\mathsf d (G) - \mathsf d^* (G)$ we refer to \cite[Theorem 1]{Ma92} and more recently to \cite{Li20a}.



\begin{theorem}\label{thm:main2}
Let $G=C_2^r \oplus C_6$ with $r\in \{5,6\}$. Then $$\nu(G)=\nu_2(G)=\dd(G)=\mathsf d^* (G)+1.$$
\end{theorem}




\section{Notation and terminology}\label{sec:criterion}

For real numbers $a,b\in \mathbb{R}$, let $[a,b]=\{x\in \mathbb{Z}: a\leq x\leq b\}$ denote the discrete interval between $a$ and $b$. Let $G$ be an additively written finite abelian group.
The identity element of $G$ is denoted by $0$. For a subset
$A\subseteq G$, the subgroup generated by $A$ is denoted by
$\langle A\rangle$. We write $C_n$ for a cyclic group of order $n$.


We use the standard language of sequences over groups;  see \cite{Ga-Ge06b}
 and \cite[Chapter 5]{Ge-HK06a}.  A sequence over
$G$ is an element of the free abelian monoid $\mathcal F(G)$ and will be written in the form
\[
S= \prod\limits_{g\in G} g^{\mathsf v_g(S)} = g_1\boldsymbol{\cdot}\ldots\boldsymbol{\cdot}g_\ell \,,
\]
where $\mathsf v_g(S)\in\mathbb N_0$ is the multiplicity of $g$ in $S$.
The {\sl length} of a sequence $S$ is
$|S|=\sum_{g\in G}\mathsf v_g(S)$.  A sequence $T \in \mathcal F(G)$ is called a {\it subsequence } of $S$ and is denoted by $T \mid S$ if  $\mathsf v_g(T) \le \mathsf v_g(S)$ for all $g\in G$.
We write
$$\sigma(S)=\sum_{i=1}^{\ell}g_i
 \quad\text{and}\quad
 \Sigma(S)=\left\{\sum_{i\in I}g_i:
                    \emptyset\ne I\subseteq[1,\ell]\right\}.$$
The sequence $S$ is zero-sum free if $0\notin\Sigma(S)$.
It is a minimal zero-sum sequence if it is nontrivial,
$\sigma(S)=0$, and every proper nontrivial subsequence
has nonzero sum.

\section{Proof of Theorems \ref{thm:main1} and \ref{thm:main2}}\label{sec:constructions}

We need the following lemmas.


\begin{lemma} \label{2.1}
Let $G = C_2^r \oplus C_{2n}$ with $n\ge3$ odd and $r\ge1$. 
Then $\mathsf d(G)=\mathsf d^*(G)$ if and only if $r\le3$.
\end{lemma}

\begin{proof}
See \cite[Theorem 3.9.7]{Ge-Gr-Zh26a}.
\end{proof}

\begin{lemma}\label{lem:davenport}
We have  $\dd(C_2^5\oplus C_6)=11$ and $\dd(C_2^6\oplus C_6)=12$.
\end{lemma}

\begin{proof}
For $r\in\{5,6\}$, Lemma~\ref{2.1} gives
$\mathsf d(C_2^r\oplus C_6)
 \ge \mathsf d^*(C_2^r\oplus C_6)+1=r+6.$
Equality holds by
\cite[p.~2]{Sc11b} (where this result is  attributed to Ponomarenko).
\end{proof}



Now we are in a position to prove Theorems~\ref{thm:main1} and \ref{thm:main2}.





\begin{proof}[Proof of Theorem~\ref{thm:main1}]
Let $G=C_2^3\oplus C_{2n}$, where $n\ge3$ is odd.
Choose a basis $(e_1,e_2,e_3,e)$ of $G$ with $\ord(e_i)=2$
for $i \in [1,3]$ and $\ord(e)=2n$. Put $a=(n+1)/2$ and set
\begin{equation}\label{eq:family}
 T_n=e^{2n-3}(e_1+ae)(e_2+ae)(e_3+ae)
 (e_1+e_2+e_3+ae).
\end{equation}
Among the last four terms, the only nontrivial subsequence whose
sum lies in $\langle e\rangle$ consists of all four terms.
Their sum is $4ae=2e$, which cannot be completed to zero by
adding at most $2n-3$ copies of $e$. A nontrivial subsequence
consisting only of copies of $e$ also has nonzero sum.
Thus $T_n$ is zero-sum free, and $|T_n|=2n+1=\mathsf d^* (G)-1=\mathsf d (G)-1$ by Lemma~\ref{2.1}.

Put $A=\{e_1+e_2,e_1+e_3,e_2+e_3\}$. For each $x\in A$,
any subsequence of $T_n$ whose sum lies in $x+\langle e\rangle$
contains exactly two of the last four terms in \eqref{eq:family}.
Its sum therefore has the form $x+(2a+t)e$, where
$0\le t\le2n-3$. Since $2a=n+1$, the coefficient $2a+t$
is congruent to neither $n-1$ nor $n$ modulo $2n$.
Consequently,
\begin{equation}\label{equation:missing-sums}
 \{x+(n-1)e,\ x+ne\}
 \subseteq G^\bullet\setminus\Sigma(T_n)
 \qquad\text{for every }x\in A.
\end{equation}

Assume to the contrary, that
$G^\bullet\setminus\Sigma(T_n)\subseteq\alpha+H$
for some subgroup $H\le G$ and some $\alpha\notin H$.
Since differences of elements of $\alpha+H$ belong to $H$,
\eqref{equation:missing-sums} gives
\begin{align*}
 e   &= (x+ne)-(x+(n-1)e)\in H,\\
 x-y &= (x+ne)-(y+ne)\in H
\end{align*}
for all $x,y\in A$. Thus $A-A\subseteq H$.
Since $A-A=A\cup\{0\}$, we obtain $A\subseteq H$.
Together with $e\in H$, this implies that for every $x\in A$,
one has $x\in H$ and hence $x+(n-1)e\in H$.
Since also $x+(n-1)e\in\alpha+H$, it follows that $x+(n-1)e\in H\cap(\alpha+H)$, contradicting $\alpha\notin H$. Since $|T_n|=\dd(G)-1$, we obtain $\nu(G)>\dd(G)-1$.
The equalities follow from \eqref{basic-inequ-2} immediately.
\end{proof}

\medskip

\begin{proof}[Proof of Theorem~\ref{thm:main2}]
By Inequality \eqref{basic-inequ-2},   it suffices to show that $\nu(G)>\dd(G)-1$. We proceed in two steps.

1. Let $G=C_2^5\oplus C_6$.
Choose a basis $(e_1,\ldots,e_5,e)$ of $G$ with
$\ord(e_i)=2$ for $i\in[1,5]$ and $\ord(e)=6$. Set
\begin{equation}\label{eq:S}
 S=\prod_{i=1}^5(e_i+e)(e_i+4e).
\end{equation}
Then $|S|=10=\dd(G)-1$ by Lemma~\ref{lem:davenport}. A subsequence whose sum lies in $\langle e\rangle$
must contain both or neither of $e_i+e$ and $e_i+4e$
for each $i\in[1,5]$. Since each pair sums to $5e=-e$,
any nontrivial such subsequence has sum $-ke$ for some
$k\in[1,5]$, which is nonzero. Thus $S$ is zero-sum free.

Put $g_0=e_1+\ldots+e_5$. A subsequence whose sum lies in
$g_0+\langle e\rangle$ contains exactly one term from each pair.
For a sum in $g_0-e_i+\langle e\rangle$, it contains exactly
one term from each pair with index different from $i$,
and both or neither of the terms in the $i$th pair.
Consequently,
\begin{align*}
 \Sigma(S)\cap(g_0+\langle e\rangle)
   &=g_0+\{2e,5e\},\\
 \Sigma(S)\cap(g_0-e_i+\langle e\rangle)
   &=g_0-e_i+\{0,e,3e,4e\}
   \qquad(i\in[1,5]).
\end{align*}
It follows that
$$
 \{g_0,g_0+e\}
 \cup\{g_0-e_i+2e:i\in[1,5]\}
 \subseteq G^\bullet\setminus\Sigma(S).
$$

Assume to the contrary, that
$G^\bullet\setminus\Sigma(S)\subseteq\alpha+H$
for some subgroup $H\le G$ and some $\alpha\notin H$.
Taking differences of the missing sums above gives
\begin{align*}
 e &= (g_0+e)-g_0\in H,\\
 -e_i+2e &= (g_0-e_i+2e)-g_0\in H
 \qquad(i\in[1,5]).
\end{align*}
Hence $e,e_1,\ldots,e_5\in H$, so $H=G$,
contradicting $\alpha\notin H$. Thus $\nu(G)>\dd(G)-1$ in this case.

\smallskip
2. Let $G=C_2^6\oplus C_6$.
Choose a basis $(e_1,\ldots,e_6,e)$ of $G$ with
$\ord(e_i)=2$ for $i\in[1,6]$ and $\ord(e)=6$.
Regard the sequence $S$ in \eqref{eq:S} as a sequence over $G$.
For any $g\in\langle e_1,\ldots,e_5,e\rangle$, set
\begin{equation}\label{eq:lift}
	S_g=S(e_6+g).
\end{equation}
Every subsequence $(e_6+g)T$ of $S_g$ with $T\mid S$ has sum
$\sigma((e_6+g)T)\in e_6+\langle e_1,\ldots,e_5,e\rangle$, which does not contain zero.
Since $S$ is zero-sum free, it follows that $S_g$ is zero-sum free
and $|S_g|=11=\dd(G)-1$ by Lemma~\ref{lem:davenport}.

Moreover, any subsequence of $S_g$ whose sum lies in
$\langle e_1,\ldots,e_5,e\rangle$ is a subsequence of $S$.
Consequently,
\[
\langle e_1,\ldots,e_5,e\rangle^\bullet\setminus\Sigma(S)
\subseteq G^\bullet\setminus\Sigma(S_g).
\]
Let $g_0=e_1+\ldots+e_5$. Thus
$$\{g_0,g_0+e\}
 \cup\{g_0-e_i+2e:i\in[1,5]\}
 \subseteq G^\bullet\setminus\Sigma(S_g).$$

Assume to the contrary that
$G^\bullet\setminus\Sigma(S_g)\subseteq\alpha+H$
for some subgroup $H\le G$ and some $\alpha\notin H$.
Since differences of elements of $\alpha+H$ belong to $H$,
we obtain
\begin{align*}
	e &= (g_0+e)-g_0\in H,\\
	e_i-2e &= g_0-(g_0-e_i+2e)\in H
	\qquad\text{for every }i\in[1,5].
\end{align*}
Together with $e\in H$, this implies that for every $i\in[1,5]$,
one has $e_i=(e_i-2e)+2e\in H$, and hence
$g_0=e_1+\ldots+e_5\in H$.
Since $g_0\in G^\bullet\setminus\Sigma(S_g)\subseteq\alpha+H$,
it follows that $g_0\in H\cap(\alpha+H)$,
contradicting $\alpha\notin H$.
Since $S_g$ is zero-sum free and $|S_g|=\dd(G)-1$, we obtain $\nu(G)>\dd(G)-1$.
\end{proof}

\begin{thebibliography}{99}
\bibitem[Em69a]{Em69a} {P. van} {Emde Boas}, \emph{A combinatorial problem on finite abelian groups {II}}, Reports ZW-1969-007, Math. Centre, Amsterdam, 1969.
\bibitem[Ga-Ge06b]{Ga-Ge06b} W.~Gao and A.~Geroldinger, \emph{Zero-sum problems in finite abelian groups: A survey}, Expositiones Mathematicae \textbf{24} (2006), 337--369. \href{https://doi.org/10.1016/j.exmath.2006.07.002} {doi:10.1016/j.exmath.2006.07.002}.
\bibitem[Ga00b]{Ga00b} W.~Gao, \emph{On Davenport's constant of finite abelian groups with rank three}, Discrete Mathematics \textbf{222} (2000), 111--124. \href{https://doi.org/10.1016/S0012-365X(00)00010-8} {doi:10.1016/S0012-365X(00)00010-8}.
\bibitem[Ge-Gr-Zh26a]{Ge-Gr-Zh26a} A.~Geroldinger and D.J.~Grynkiewicz and Q.~Zhong, \emph{Combinatorial Factorization Theory. Factorization Theory of Krull Monoids via Additive Combinatorics}, Mathematical Surveys and Monographs, 296. American Mathematical Society, Providence, RI, (2026), \href{https://bookstore.ams.org/view?ProductCode=SURV/296}{doi: 10.1090/surv/296}.
\bibitem[Ge-HK06a]{Ge-HK06a} A. Geroldinger and F. Halter-Koch, \emph{Non-Unique Factorizations: Algebraic, Combinatorial and Analytic Theory}, Pure and Applied Mathematics 278, Chapman \& Hall/CRC, 2006.
\bibitem[Ge-Ya27a]{Ge-Ya27a} A.~Geroldinger and W.~Yang, \emph{On a classical zero-sum invariant}, \href{https://arxiv.org/abs/2608.19090v1}{arXiv:2608.19090v1} [math.NT], 19 August 2026.
\bibitem[Li20a]{Li20a} C.~Liu, \emph{On the lower bounds of {D}avenport constant}, J. Comb. Theory, Ser. A {\bf 171} (2020), {Paper No. 105162, 15pp}.
\bibitem[Ma92]{Ma92} M.~Mazur, \emph{A note on the growth of {D}avenport's constant}, Manuscr. Math. {\bf 74} (1992), 229 -- 235. \href{https://doi.org/10.1007/BF02567669}{doi:10.1007/BF02567669}.
\bibitem[Sc11b]{Sc11b} W.~A.~Schmid, \emph{The inverse problem associated to the Davenport constant for $C_2\oplus C_2\oplus C_{2n}$, and applications to the arithmetical characterization of class groups}, The Electronic Journal of Combinatorics \textbf{18}(1) (2011), Paper P33. \href{https://doi.org/10.37236/520}{doi:10.37236/520}.
\end{thebibliography}
\end{document}
